\section{Discussion}\label{sec:discussion}
% statements-of-record rows resolved here: none

The standard BSD formula is a product of scalar factors: an analytic
leading coefficient divided by a period, a regulator, and a finite factor
built from Tamagawa numbers, torsion and $\Sha$, with local $p$-adic
normalizations and determinant conventions hidden in the assembly.  The
formula is a powerful summary, but each factor is the image of richer
data.  The question of this paper is which of those data a given readout
retains.

The Schur residual $\XiC$ turns that question into finite linear algebra.
If all native coordinates of a column are admitted, the residual vanishes.
In general, for the declared target, the residual measures the variation
of the target that the admitted source coordinates do not capture; a
source that discards some coordinates can still be adequate for a target
that ignores them, as the orientation-normalized readout of
\Cref{thm:apparatus:five-col-bk-comparison} shows.  Most of the adequacy
collapses in this paper are elementary for exactly this reason.  Their
value is not depth but bookkeeping: they say which coordinates must be
present before a datum can serve as a source for a given readout.

The ``does not determine'' results complement the collapses.  Some are
formal (\Cref{thm:apparatus:finite-scalar-public-shadow,thm:apparatus:gram-matrix-shadow,thm:apparatus:finite-window-no-go,thm:apparatus:finite-prime-cover-no-go}),
and their content is the precise identification of what is forgotten.
One is arithmetic: the curves $571a1$ and $2045b1$
(\Cref{ex:apparatus:sha-pair}) have the same rank, the same Tamagawa
numbers, the same torsion and the same $2$-torsion dimension of $\Sha$,
but $2$-primary parts of $\Sha$ of orders $4$ and $16$.  The dimension of
$\Sha[2]$ is the datum most directly computed by $2$-descent, and the
example shows concretely that it is not the datum the BSD formula needs.

Three normalization points recur and deserve emphasis.  First, a
Cassels--Tate pairing with values in $\Q/\Z$ does not give a canonical
numerical Pfaffian.  The invariant finite factor is the order of the
paired group; an alternating integral presentation $A$ gives
$\#S_p=|\det A|=\Pf(A)^2$, which determines $|\Pf(A)|$ but not its sign
(\Cref{sec:finite-source}).  Second, on the finite
local condition the relevant map is the logarithm; the dual exponential
vanishes there (\Cref{sec:p-adic}).  Third, the orientation enters the
determinant assembly as an independent input; the assembled value must
not be reused as a factor (\Cref{sec:det-assembly,sec:comparison}).  Each
of these is a small point, but each changes a formula.

At the two ends of the apparatus sit two calibrations.  At the source end,
the exterior algebra $\vsrc_r$ records the shape that a rank-$r$ source
must have, and the rank-two discussion shows that a scalar leading-term
identity does not by itself provide such a source.  At the target end,
the normalization $\kapr=\#\Sha/\#E(\Q)_{\tors}^{\,2}$ is exactly what
makes the cascade form of the leading term coincide with the standard
one.

The language of typed columns, readouts and residuals comes from the Six
Birds framework~\citep{TsiokosNeedles2026,TsiokosFoundationsII2026,TsiokosFoundationsIII2026},
but nothing in the paper depends on that framework beyond the definitions
given here.  The apparatus is used by the companion
paper~\citep{TsiokosBSDClosure2026}, which treats the scalar BSD identity
as a closure problem under explicit hypotheses.
